% Created 2015-11-18 Wed 09:32
\documentclass[bigger,allowframebreaks]{beamer}
\usepackage[utf8]{inputenc}
\usepackage[T1]{fontenc}
\usepackage{fixltx2e}
\usepackage{graphicx}
\usepackage{longtable}
\usepackage{float}
\usepackage{wrapfig}
\usepackage{rotating}
\usepackage[normalem]{ulem}
\usepackage{amsmath}
\usepackage{textcomp}
\usepackage{marvosym}
\usepackage{wasysym}
\usepackage{amssymb}
\usepackage{hyperref}
\tolerance=1000
\mode<beamer>{\useinnertheme{rounded}\usecolortheme{rose}\usecolortheme{dolphin}\setbeamertemplate{navigation symbols}{}\setbeamertemplate{footline}[frame number]{}}
\mode<beamer>{\usepackage{amsmath}\usepackage{ae,aecompl,sgamevar}}
\let\oldframe\frame\renewcommand\frame[1][allowframebreaks]{\oldframe[#1]}
\setbeamertemplate{frametitle continuation}[from second]
\usetheme{default}
\author{Christoph Schottmüller}
\date{}
\title{Game Theory: Global Games}
\hypersetup{
  pdfkeywords={},
  pdfsubject={},
  pdfcreator={Emacs 24.5.1 (Org mode 8.2.10)}}
\begin{document}

\maketitle
\begin{frame}{Outline}
\tableofcontents
\end{frame}




\section{Global Games: Stag Hunt}
\label{sec-1}
\begin{frame}[label=sec-1-1]{Stag Hunt}
\begin{exampleblock}{Example}
  \begin{figure}[h]
    \centering
   \begin{game}{2}{2}
          \>  H2 \> S2\\
     H1   \> 3,3 \> 3,0\\
    S1   \> 0,3 \> 4,4
   \end{game}
  \caption{Stag hunt}
\end{figure}

\begin{itemize}
\item two pure NE (+ 1 mixed NE) $\Rightarrow$ which should we predict as outcome of the game?
\end{itemize}
\end{exampleblock}
\end{frame}

\begin{frame}[label=sec-1-2]{Two approaches}
\begin{itemize}
\item payoff dominance
\item risk dominance
\begin{itemize}
\item multiply the deviation losses of both players in a given equilibrium and select the equilibrium where the product is highest ($here: (3-0)^2>(4-3)^2$)
\item intuitively: which equilibrium is supported by a larger set of beliefs? i.e. if the other player mixes $(1/2,1/2)$, what do you choose?
\end{itemize}
\end{itemize}
\end{frame}

\begin{frame}[allowframebreaks,label=]{Global Games: Carlsson and van Damme (1993)}

\begin{figure}[h]
 \centering
 \begin{game}{2}{2}
 \> H2 \> S2\\
 H1   \>$x,x$\>$x,0$\\
 S1 \>$0,x$\>$4,4$
\end{game}
\end{figure}

\begin{itemize}
\item state of the world $x\in[-1,5]$ is unknown to players
\item assume each state is equally likely ($x$ distributed uniformly on $[-1,5]$)
\item both players observe a private signal $x_i\in[x-\varepsilon ,x+\varepsilon ]$ for some $\varepsilon >0$
\item assume the two $x_i$ are independently drawn from a uniform distribution on $[x-\varepsilon ,x+\varepsilon ]$
\item interpretation: players are not exactly sure which game they play
\end{itemize}

\framebreak

assume $\varepsilon=0.1$
\begin{itemize}
\item you get a signal $x_1=-0.5$
\begin{itemize}
\item what do you believe about the state of the world?
\item what do you believe about the other player's signal?
\item which action should you take?
\end{itemize}

\item you get signal $x_1<0$
\begin{itemize}
\item \dots{}
\end{itemize}

\item you get a signal $x_1=0$
\begin{itemize}
\item \dots{}
\end{itemize}

\item what is the biggest $x_1$ for which you play $S1$ for sure?
\end{itemize}

\framebreak
\begin{itemize}
\item you get signal $x_1>4$
\begin{itemize}
\item what do you believe about the state of the world?
\item what do you believe about the other player's signal?
\item which action should you take?
\end{itemize}

\item what is the smallest $x_1$ for which you play $H1$ for sure?
\end{itemize}

\framebreak

\begin{itemize}
\item result does not depend on $\varepsilon =0.1$ but is true for any $\varepsilon >0$ that is not too big
\item if $\varepsilon \rightarrow 0$, we get a unique prediction for any stag hunt game
\item Carlsson and van Damme show that the global game approach selects the risk dominant equilibrium in $2\times2$ games with two strict equilibria
\item result does not depend on uniform distributions
\end{itemize}
\end{frame}

\begin{frame}[allowframebreaks,label=]{Intermezzo: Cutoff Strategies}

\begin{itemize}
\item A strategy in the global game is a function that associates an action with each signal
\end{itemize}
$$s_i(x_i): \;[-1,5]\rightarrow\{Hi,Si\}$$
\begin{itemize}
\item we define a \emph{cutoff strategy} with cutoff $\bar s$ as
\end{itemize}
$$ s_i(x_i)=\begin{cases}
Hi & \text{ if }x_i\geq \bar s_i\\
Si & \text{ if }x_i<\bar s_i
\end{cases}$$

\framebreak

if $i$ plays a cutoff strategy with cutoff $\bar s_i\in(-\varepsilon,4+\varepsilon)$ which cutoff strategy is a best response of $j$ (i.e. what is the best response cutoff $\bar s_j$)?
\begin{itemize}
\item if $j$ gets the signal $x_j=\bar s_j$, he should be indifferent between Hj and Sj
\item can $j$ be indifferent if $\bar s_i\leq \bar s_j-2\varepsilon$?
\item can $j$ be indifferent if $\bar s_i\geq \bar s_j+2\varepsilon$?
\item after observing $x_j$, what does $j$ believe about the distribution of $x_i$?
\end{itemize}
\begin{itemize}
\item can $\bar s_j=\bar s_i$ be a best response? what would $j$ think about the probability that $i$ plays Si when receiving signal $x_j=\bar s_j$ in this case?
\end{itemize}
\begin{itemize}
\item if $\bar s_i<2$, (i) should $\bar s_j$ be above or below $\bar s_i$ to make $j$ indifferent? (ii) could $\bar s_j>2$ make $j$ indifferent?
\end{itemize}
\end{frame}
\begin{frame}[label=sec-1-5]{Global Games and iterative elimination of strictly dominated strategies}
\begin{itemize}
\item back to the global game: which strategies are strictly dominated?
\end{itemize}
\begin{itemize}
\item given that the other player does not play strictly dominated strategies, which of my strategies are strictly dominated?
\item iterative elimination of dominated strategies
\end{itemize}
\begin{block}{Result}
Iterative elimination of strictly dominated strategies leads to a unique prediction in the global game.
\end{block}
\end{frame}


\begin{frame}[label=sec-1-6]{Global Games and Knowledge}
\begin{itemize}
\item if $i$ receives signal $x_i$, what does he know about
\begin{itemize}
\item state of the world
\item which states $j$ considers possible
\item which states $j$ thinks $i$ considers possible
\item \dots{}
\end{itemize}
\end{itemize}
\begin{itemize}
\item global games capture -- in a tractable way -- the idea that 
\begin{itemize}
\item payoffs are not common knowledge
\item beliefs are not common knowledge
\end{itemize}
\end{itemize}
\end{frame}



\section{An investment example}
\label{sec-2}
\begin{frame}[allowframebreaks,label=]{An investment example}
\begin{itemize}
\item Coordination problems as above are relevant in many settings
\item 2 investors
\item investors decide (simultaneously) whether to invest in a project or not ($A_i=\{I,N\}$)
\item investing carries a cost of 1
\item the project succeeds only if both invest
\item a successfull project generates revenue of 4 which is equally split
\item not investing in the project gives a payoff of $r$
\end{itemize}
\begin{figure}[h]
       \centering
      \begin{game}{2}{2}
             \>  I \> N\\
       I  \> $1,1$ \> $-1,r$\\
       N   \> $r,-1$ \> $r,r$
      \end{game}
     \caption{Investment game}
   \end{figure}

\begin{itemize}
\item for which values of $r$ does an investor have a dominant strategy?
\item say, both investors have the prior that $r$ is distributed uniformly on $[-2,2]$ and receive a private signal $x_i$ which is uniformly distributed on $[r-0.1,r+0.1]$
\item what is the belief of $i$ after observing $x_i$ concerning 
\begin{itemize}
\item the state $r$
\item the signal of the other investor?
\end{itemize}
\item for which values of $x_i$ is $I$ strictly dominated by $N$?
\item given that the other investor does not use dominated strategies, why is $I$ strictly dominated by $N$ for values of $x_i$ slightly below 1?
\item What is the equilibrium (in cutoff strategies)?
\end{itemize}
\end{frame}


\section{Revision questions and exercises}
\label{sec-3}
\begin{frame}[label=sec-3-1]{Review questions}
\begin{itemize}
\item What is the Stag Hunt game and what is the basic problem described by this game?
\item What is the main game theoretic challenge in the Stag Hunt?
\item What is a global game?
\item What does a player know in a global game? What is common knowledge?
\item How does the global game approach select an equilibrium?
\item Why is it surprising that the global game approach selects a unique equilibrium?
\item reading: Carlsson and van Damme, ECTA 93, p. 989-993; *Morris and Shin, 2000, p. 1-15; *Morris and Shin, 1998
\end{itemize}
\end{frame}

\begin{frame}[allowframebreaks,label=]{Exercises}

\begin{itemize}
\item Compare equilibrium selection by the global game approach with equilibrium refinement by perfect equilibrium.
\end{itemize}
\begin{itemize}
\item This exercise is about a model of speculating against a fixed exchange rate: Two players have to decide whether they want to speculate against the peg or not. If both players speculate against, the central bank is unable to defend the peg and the speculators make a profit of $x$. If only one player speculates against the peg, the central bank will successfully defend the peg and the speculating player receives $x-1$. This gives payoffs as below
\end{itemize}
\begin{figure}[h]
       \centering
      \begin{game}{2}{2}
             \>  S \> N\\
        S  \> $x,x$ \> $x-1,0$\\
       N   \> $0,x-1$ \> $0,0$
      \end{game}
     \caption{Speculating against a currency peg}
   \end{figure}
Each player does not know $x$ but receives a signal $x_i$ which is uniformly distributed on $[x-0.1,x+0.1]$ (the two signals are independent).
\begin{itemize}
\item For which values of $x$ is S a dominant action?
\item For which values of $x$ is N a dominant action?
\item Suppose player 1 observes $x_1=0.9$. What does he believe about the distribution of $x$ (assume that players had a uniform prior on $[-1,2]$ about the true state $x$ before observing $x_i$)? What does he believe about the distribution of $x_2$?
\end{itemize}
\begin{itemize}
\item Suppose player 1 believes that player 2 plays $S$ whenever $x_2\geq 1$. What should player 1 do when observing signal $x_1=0.9$?
\end{itemize}
\begin{itemize}
\item Suppose player 1 believes that player 2 plays S whenever $x_2\geq \bar x_2$ (and let $\bar x_2\geq 1/2$). For which values of $x_1$ should player 1 definitely play S himself? *What is 1's best response?
\end{itemize}
\begin{itemize}
\item Can you find a symmetric Nash equilibrium in this game? (recall that a strategy is a function assigning an action to each signal)
\end{itemize}
\begin{itemize}
\item *Show that only cutoff strategies are optimal and that the equilibrium of the previous subquestion is the unique Nash equilibrium. Show that this NE is the unique strategy surviving iterative elimination of strictly dominated strategies. (hint: check the Morris/Shin paper for a solution)
\end{itemize}
\end{frame}
% Emacs 24.5.1 (Org mode 8.2.10)
\end{document}